\chapter{Research and Existing Solutions}
\label{chap:research}

\section{Research Method and Literature Survey Approach}
To establish a rigorous theoretical foundation and address the FYP Committee's mandate regarding state-of-the-art Dynamic Difficulty Adjustment (DDA) and Procedural Content Generation via Reinforcement Learning (PCGRL), the project team conducted a structured narrative literature review (not a formal systematic review: no database search protocol or screening counts were recorded) covering \textbf{40 sources} (peer-reviewed publications, preprints, foundational texts and a small number of industry sources) across game AI, human-computer interaction, and machine learning venues (IEEE Transactions on Games, IEEE Conference on Games, ACM CHI, Foundations of Digital Games, and Springer Multimedia Tools and Applications).

Our investigation followed a structured three-tier review methodology:
\begin{enumerate}
    \item \textbf{Generational DDA Survey:} Examining the evolution of difficulty adjustment algorithms from heuristic threshold triggers to modern statistical, continuous, and personalized reinforcement learning models~\cite{mortazavi2024dda, zohaib2018dynamic, fisher2024exploring, pfau2020enemy, fuchs2024personalized, rahim2023continuous, romeo2025ntrl, xue2017dynamic}.
    \item \textbf{PCGRL State-of-the-Art Analysis:} Reviewing learned procedural generators conditioned on target features, open-ended environments, multi-objective Pareto optimization, GPU-accelerated environments, graph topologies, and experience-driven, player-adaptive generation~\cite{khalifa2020pcgrl, earle2021illuminating, earle2021controllable, yuan2026mopcgrl, li2025measuring, bazzaz2026imitation, earle2024pcgrlplus, rupp2024gpcgrl, yannakakis2018ai, yannakakis2011experience, shaker2010personalized, shu2021edrl, wang2022oparl}.
    \item \textbf{Human-Centered Psychometric Instruments:} Evaluating standardized psychological scales designed to measure player flow, immersion, perceived competence, and autonomy in real-time gaming environments~\cite{csikszentmihalyi1990flow, denisova2016convergence, tyack2020sdt, michailidis2018flow}.
\end{enumerate}

\section{Existing Solutions and Products}
The academic and commercial landscape of adaptive content generation is grouped below into three paradigms. They are conceptual families rather than strict eras: rule-based systems remain common in recent commercial games, and some statistical methods are recent.

\subsection{Generation 1: Heuristic and Rule-Based Systems}
Early commercial difficulty adaptation relied on deterministic conditional triggers ($\text{IF}/\text{THEN}$ rules) operating on scalar difficulty ranks. Notable examples include:
\begin{itemize}
    \item \textbf{Resident Evil 4 (Capcom):} Is widely reported to use a hidden difficulty rank that rises when the player progresses quickly and cleanly, making enemies hit harder and act more aggressively~\cite{engadget2015re4}. Capcom has not officially documented the mechanism, so its exact inputs and effects are uncertain.
    \item \textbf{Left 4 Dead (Valve):} Features the ``AI Director''~\cite{booth2009l4d}, which estimates each survivor's emotional intensity (rising with damage taken and nearby enemy deaths) and cycles population pacing through build-up, peak and relax phases.
    \item \textbf{Hades (Supergiant Games):} Offers an optional, player-enabled ``God Mode'' providing $20\%$ damage resistance that grows by $2\%$ with each death, up to $80\%$. It is player-chosen assistance rather than automatic DDA, but shows the same scalar-multiplier design.
\end{itemize}
\textbf{Limitation:} These systems operate via scalar multipliers, suffer from high transparency, and frequently trigger the ``rubber-banding penalty,'' wherein players feel patronized upon discovering that their struggles artificially weaken enemies.

\subsection{Generation 2: Statistical and Probabilistic Modeling}
To mitigate the rigidity of rule-based triggers, second-generation systems incorporated probabilistic inference, Markov Decision Processes (MDPs), and fuzzy logic controllers:
\begin{itemize}
    \item \textbf{Fuzzy-Logic Difficulty Controllers:} Chrysafiadi et al.~\cite{chrysafiadi2023fuzzy} designed a fuzzy-logic controller for an educational 3D game that maps imprecise descriptors of player performance onto graded difficulty adjustments.
    \item \textbf{Engagement Optimization (Xue et al.~\cite{xue2017dynamic}):} Models player progression as a probabilistic graph and selects difficulty to maximize expected engagement (rounds played before churn) in a commercial game.
\end{itemize}
\textbf{Limitation:} These methods are smoother than Generation 1 but usually adjust one or a few scalar difficulty levels, and they are hand-designed per game.

\subsection{Generation 3: Deep Reinforcement Learning}
Modern research frames game generation and difficulty adaptation as reinforcement learning problems:
\begin{itemize}
    \item \textbf{PCGRL Benchmark (Khalifa et al.~\cite{khalifa2020pcgrl}):} Formulates level generation as an MDP where an agent acts as an iterative brush placing tiles to maximize path length and reachability.
    \item \textbf{Controllable Content Generators (Earle et al.~\cite{earle2021controllable}):} Conditions PCGRL generators on designer-specified numerical control targets (e.g., target path length).
    \item \textbf{Quality-Diversity Generators and Multi-Objective PCGRL:} Earle et al.~\cite{earle2021illuminating} evolve a diverse archive of neural-cellular-automata level generators with quality-diversity search (an evolutionary alternative to RL), and Yuan et al.~\cite{yuan2026mopcgrl} extend PCGRL to multiple simultaneous objectives (Pareto trade-offs), which is relevant to FORGE's competing goals of challenge, solvability and frame budget. Li et al.~\cite{li2025measuring} survey metrics for measuring the diversity of generated scenarios, and Bazzaz and Cooper~\cite{bazzaz2026imitation} show that players' enjoyment ratings of generated levels depend on whether they believe a human or an AI created them.
    \item \textbf{PCGRL+ Scaled Level Generators (Earle et al.~\cite{earle2024pcgrlplus}):} Demonstrates billion-step parallel level generation on GPUs, overcoming long-standing training bottlenecks in reinforcement learning procedural level generators.
    \item \textbf{Graph-Based PCGRL (Rupp \& Eckert~\cite{rupp2024gpcgrl}):} Extends PCGRL beyond 2D tile matrices to graph data (e.g., game economies and skill trees) by editing a graph's adjacency matrix, a representation that could in principle also describe room adjacency and progression graphs.
    \item \textbf{Experience-Driven and Player-Adaptive PCGRL (Shu et al.~\cite{shu2021edrl}; Wang \& Liu~\cite{wang2022oparl}):} Generate \textit{Super Mario Bros.} level segments online with RL, rewarding segments that suit a model of the player's experience; OPARL further adapts generated difficulty to artificial players' play speed in real time. Together with earlier player-model-driven personalization~\cite{shaker2010personalized}, these are the closest prior work to FORGE. They differ from FORGE in genre (platformer level geometry rather than roguelike combat), in the adapted quantities, and in being evaluated outside a shipped game engine.
    \item \textbf{Encounter Generation for DDA (Romeo \& Bagdanov~\cite{romeo2025ntrl}):} Proposes NTRL (2025), which generates tabletop role-playing combat encounters from the real-time attributes of the party, framed as a contextual bandit, to hit a high (about 70\%) win-rate target.
    \item \textbf{Personalized Imitation-RL DDA (Fuchs et al.~\cite{fuchs2024personalized}):} Trains one agent to imitate the individual player and a second agent to defeat that imitation, personalizing opponent difficulty in a fighting-game setting.
    \item \textbf{Continuous Difficulty-Space DDA (Rahim et al.~\cite{rahim2023continuous}):} Replaces the small discrete state--action spaces of earlier RL-based DDA with a continuous difficulty metric in a visual working-memory game, reporting better competence ratings than two rule-based baselines in a 52-participant study, and a smaller score decline over a session than one of them.
    \item \textbf{Long-Term Adaptation Effects (Pfau et al.~\cite{pfau2020enemy}):} Reports a long-term study in a live MMORPG in which deep player-behavior models for DDA improved intrinsic motivation compared with heuristic parameter tuning, whose adjustments players often notice (e.g., rubber-banding).
    \item \textbf{Industrial Engagement Optimization (Xue et al.~\cite{xue2017dynamic}):} Formulates DDA for a commercial game as an optimization problem that selects difficulty to maximize player engagement (expected rounds played before churn).
    \item \textbf{DDA Methods Comparison (Fisher \& Kulshreshth~\cite{fisher2024exploring}):} Compares five difficulty conditions (fixed, user-selected, performance-based, emotion-based and hybrid) in a custom first-person shooter within-subjects study, finding no single method that clearly beat static difficulty, which underlines the difficulty of tuning adaptation well.
\end{itemize}

\section{Comparison of Existing Solutions}
Table~\ref{tab:lit_comparison} evaluates Project FORGE against representative systems from both commercial gaming and peer-reviewed literature.

\begin{table}[!htbp]
\centering
\caption{Comparison of Game Difficulty Adjustment and PCG Paradigms}
\label{tab:lit_comparison}
\scriptsize
\begin{tabular}{p{0.20\textwidth}p{0.14\textwidth}p{0.14\textwidth}p{0.13\textwidth}p{0.12\textwidth}p{0.13\textwidth}}
\toprule
\textbf{System / Reference} & \textbf{Paradigm} & \textbf{Adaptation Loop} & \textbf{Action Space} & \textbf{Solvability} & \textbf{In-Engine Overhead} \\
\midrule
\textit{Resident Evil 4}~\cite{engadget2015re4} & Gen 1 Heuristic & Reactive scalar rank (unofficial) & Enemy damage/aggression & Fixed design & Negligible \\
\textit{Left 4 Dead} AI Director~\cite{booth2009l4d} & Gen 1 Heuristic & Pacing state machine & Spawn wave timer & Fixed design & Low \\
Fuzzy DDA~\cite{chrysafiadi2023fuzzy} & Gen 2 Statistical & Reactive fuzzy rules & Graded multipliers & Fixed design & Low \\
\texttt{gym-pcgrl}~\cite{khalifa2020pcgrl} & Gen 3 RL & Open-loop target & Discrete tile grid & Heuristic reward & High (Python env) \\
Controllable PCGRL~\cite{earle2021controllable} & Gen 3 RL & Open-loop designer & Scalar target vector & Path-length reward & High (External) \\
NTRL~\cite{romeo2025ntrl} & Gen 3 RL & Party-state bandit & Encounter composition & N/A (tabletop) & Not reported \\
EDRL / OPARL~\cite{shu2021edrl, wang2022oparl} & Gen 3 RL & Online, player-adaptive & Next level segment & Playability-checked segments & Python (research setting) \\
\textbf{FORGE (Ours)} & \textbf{Gen 3 RL + DDA} & \textbf{Closed-loop telemetry} & \textbf{Continuous 8D (7 settings + archetype mix)} & \textbf{Connected by construction (BFS-checked)} & \textbf{$<1$\,ms per floor (GDScript)} \\
\bottomrule
\end{tabular}
\end{table}

\section{Limitations and Gaps in Existing Literature}
The review points to three gaps that FORGE targets:
\begin{enumerate}
    \item \textbf{Limited Closed-Loop Personalization for Combat Content:} Most PCGRL generators are conditioned on random seeds or designer-set targets~\cite{earle2021controllable, earle2024pcgrlplus} rather than on a model of the individual player's experience, as advocated by experience-driven PCG~\cite{yannakakis2011experience}. Player-adaptive RL generators exist~\cite{shu2021edrl, wang2022oparl}, as does player-model-driven personalization~\cite{shaker2010personalized}, but they address platformer level geometry and one or two adapted features, with simulated or artificial players. Closed-loop adaptation of several combat and environment attributes in an action roguelike, driven by in-game telemetry, remains comparatively unexplored. Large Language Models are being explored for high-level game design~\cite{gallotta2024llmsurvey}, but their latency and unpredictability make them ill-suited to per-floor runtime combat modulation.
    \item \textbf{The Scalar Combat Balancing Oversimplification:} Most DDA systems reduce difficulty to a one-dimensional multiplier (e.g., enemy health $+10\%$). Combat difficulty is arguably combinatorial: a high-speed adversary with low health requires quick reflexes, whereas a slow, heavily armored adversary requires spatial kiting. Few existing DDA systems adjust such attributes jointly. (Section~\ref{sec:policy_eval} reports how far the current FORGE policy achieves this.)
    \item \textbf{Deployment Overhead:} Research PCGRL systems typically run inside Python and PyTorch. Engine-side inference is possible (e.g., ONNX models in Godot RL Agents~\cite{godotrl}), but it adds a native or C\# runtime dependency. FORGE instead uses a dependency-free GDScript forward pass, which is feasible because its policy network is small and queried only once per floor.
\end{enumerate}

\section{Proposed Solution: The FORGE Architecture}\label{sec:proposed}
FORGE directly addresses these limitations through a \textbf{Two-Stage Decoupled System Architecture}:
\begin{itemize}
    \item \textbf{Stage 1: Deterministic Spatial Topology Synthesis.} The dungeon structure is generated by a random walk with branching on a grid, where each new room is attached to an existing room through a matching door pair. The resulting room graph is connected by construction, so the exit room is always reachable from the start room; the CI harness checks this with a breadth-first search over the door graph on 500 seeds for each room count from 3 to 9.
    \item \textbf{Stage 2: Closed-Loop PPO Combat Ecology.} A Proximal Policy Optimization agent observes a normalized 8-dimensional state vector ($s_t \in \mathbb{R}^8$: three player-performance signals, four current difficulty settings and run depth) and outputs a continuous 8D action vector ($a_t \in [-1, 1]^8$). This action sets enemy attributes (health, damage, speed, aggro radius), hazard density, resource scarcity, room count and the mix of five enemy archetypes.
    \item \textbf{Zero-Dependency Embedded Inference:} The trained policy weights are exported from PyTorch to a compact JSON file. At runtime, a matrix multiplication routine in native GDScript executes the neural forward pass in well under 1\,ms (measured), so the game runs wherever GDScript runs, with no external dependencies.
\end{itemize}

\section{Complex Computing Problem (CCP) Justification}
Under the FAST-NUCES and Washington Accord engineering criteria, Project FORGE qualifies as a \textbf{Complex Computing Problem (CCP)} across multiple dimensions:
\begin{itemize}
    \item \textbf{WP-1 (In-depth Computing Knowledge):} Integrates advanced reinforcement learning (PPO clipped surrogate optimization), state-space observation modeling, procedural level synthesis algorithms, and real-time game engine physics.
    \item \textbf{WP-2 (Conflicting Requirements \& Multi-Objective Trade-offs):} The generator must simultaneously balance competing objectives: maintaining player challenge within the flow channel, guaranteeing topological solvability, avoiding repetitive level layouts, and maintaining a strict $16.6$\,ms (60\,FPS) rendering frame budget.
    \item \textbf{WP-3 (Non-Trivial Mathematical \& Analytical Formulation):} Involves non-linear observation-action mappings, non-uniform difficulty-weighted temporal aggregation ($\sum w_i S_i$), matrix algebra for forward propagation in GDScript, and clipped policy-gradient optimization.
    \item \textbf{WP-4 (Infrequent \& Novel Problems):} Transfers a policy trained against a simulated player (sim-to-real) into a lightweight client-side game engine, without external Python/ONNX runtimes, where it must act on real players' telemetry.
    \item \textbf{WP-5 (Engineering Integration):} Bridges offline deep RL training frameworks (Stable-Baselines3, PyTorch, Gymnasium) with real-time 2D game architecture (Godot 4.7).
    \item \textbf{WP-6 (Significance of Non-Functional Constraints):} Strict latency bounds ($<1.0$\,ms inference), zero-crash reliability under missing weights, and client-side privacy governance.
\end{itemize}

\section{Technical Challenges Anticipated and Encountered}
During design and initial implementation, the team encountered and resolved three major technical hurdles:
\begin{enumerate}
    \item \textbf{Sim-to-Real Transfer and Player Modeling:} In offline training, real human playtesters cannot provide the millions of frames required by PPO. We engineered a stochastic simulated player whose skill drifts smoothly and whose clearing time is a function of room difficulty, drawing on automated playtesting for parameter tuning, curiosity-conditioned trajectories, and procedural personas~\cite{holmgard2019procedural, zook2014playtesting, sestini2024automated}.
    \item \textbf{Reward Hacking During Offline Policy Training:} In the first version of the reward, completion time was compared with a fixed total-time target. Because total time equals room count times time per room, the PPO agent could satisfy the target for fast players by adding rooms instead of making combat harder (Section~\ref{sec:reward_history}). We addressed this by judging pace per room and adding a penalty for room counts that do not track the combat challenge (Equation~\ref{eq:reward_function}). At Mid-1 the policy still exploited the simulator in another way (enemy damage at its minimum and speed at its maximum for everyone); simulators v2 and v3 fixed this, and the v3 policy now uses all eight outputs (Section~\ref{sec:policy_eval}).
    \item \textbf{Difficulty Thrashing and Player Frustration:} Sudden spikes in player clearing time could cause an unconstrained policy to swing wildly between extreme difficulty tiers. We mitigate this with PPO's clipped objective during training, hard bounds on every output parameter (Table~\ref{tab:param_ranges}) and the cold-start anchor. The runtime has no explicit inter-floor rate limiter; with the v3 policy the mean change per floor is only $0.058$ in action units (range $-1$ to $1$) (Section~\ref{sec:shift_test}), so a limiter is considered for FYP-2.
\end{enumerate}

\section{Research Findings and Core Design Decisions}

\subsection{Why PPO Over SAC, DQN, and Contextual Bandits?}
A core inquiry raised during the Proposal Defense was the theoretical justification for selecting PPO over alternative reinforcement learning paradigms. Table~\ref{tab:rl_comparison} details this algorithmic trade-off.

\begin{table}[!htbp]
\centering
\caption{Algorithmic Comparison: Justifying PPO for Dynamic Difficulty Adjustment}
\label{tab:rl_comparison}
\small
\begin{tabular}{p{0.18\textwidth}p{0.2\textwidth}p{0.27\textwidth}p{0.25\textwidth}}
\toprule
\textbf{Algorithm} & \textbf{Action Space} & \textbf{Sample / Stability Profile} & \textbf{Suitability for DDA} \\
\midrule
\textbf{DQN}~\cite{mnih2015human} & Discrete Only & Off-policy with replay; prone to Q-value overestimation & Poor: discretizing 8 continuous settings causes a combinatorial explosion. \\
\textbf{Contextual Bandits} & Discrete/Continuous & Myopic single-step; cannot model multi-room credit assignment & Poor: Ignores long-term fatigue across multi-floor dungeon runs. \\
\textbf{SAC}~\cite{haarnoja2018soft} & Continuous & Off-policy; sample efficient; entropy-regularized exploration & Viable alternative: would also suit this action space; not chosen mainly for implementation simplicity. \\
\textbf{PPO (Ours)}~\cite{schulman2017proximal} & Continuous & On-policy; clipped surrogate keeps updates small in practice (no formal monotonic-improvement guarantee) & \textbf{Good:} stable, well-supported baseline (Stable-Baselines3); robust to noisy simulated episodes. \\
\bottomrule
\end{tabular}
\end{table}

\subsection{Rationale: PPO's Clipped Objective and Training Stability}\label{sec:ppo_proof}
Difficulty adjustment must cope with sudden player skill shocks (e.g., an expert dying to an accidental trap trigger, or a novice speedrunning a sparsely populated room). During training, such outlier episodes produce large advantage estimates, and an unconstrained policy-gradient step can then change the policy too much at once.

PPO reduces this risk by optimizing the clipped surrogate objective:
\begin{equation}
L^{CLIP}(\theta) = \hat{\mathbb{E}}_t \left[ \min\left(r_t(\theta)\hat{A}_t, \, \text{clip}\left(r_t(\theta), \, 1-\epsilon, \, 1+\epsilon\right)\hat{A}_t\right) \right]
\label{eq:ppo_clip}
\end{equation}
where $r_t(\theta) = \frac{\pi_\theta(a_t | s_t)}{\pi_{\theta_{old}}(a_t | s_t)}$ represents the probability ratio of the proposed action under the updated versus previous policy, $\hat{A}_t$ is the estimated advantage function, and $\epsilon \approx 0.2$ is the clipping boundary.

Once $r_t(\theta)$ leaves $[1-\epsilon, 1+\epsilon]$ in the direction favoured by $\hat{A}_t$, the clipped term stops contributing gradient, so the objective gives no incentive to push the policy further in that update. This keeps \emph{training} updates small in practice and makes learning more robust to noisy or extreme episodes. It is a heuristic rather than a guarantee: clipping does not strictly bound the ratio, does not enforce a formal trust region~\cite{wang2019truly}, and PPO carries no monotonic-improvement guarantee. At deployment the actor is a fixed network, so how smoothly difficulty evolves across floors depends on the learned mapping and on the hard output bounds; we therefore measure it (Section~\ref{sec:policy_eval}) instead of assuming it.

\subsection{Preserving the Principle of Change Blindness}
In perceptual psychology, ``change blindness'' refers to people's failure to notice changes to a visual scene, including gradual ones~\cite{simons2005change}; we use it here as an analogy for difficulty changes, not as an established result about them. The design goal is that difficulty parameters ($\vec{E}, \vec{M}$) change gradually between floors so that players do not notice the intervention; the current prototype bounds them to fixed ranges, and an explicit per-floor rate limiter that enforces gradual change is scheduled for FYP-2. The aim is for players to experience the game as fair and consistently challenging, supporting their sense of competence~\cite{tyack2020sdt}.
